\item {\bf Building Intuition for Bag-of-Words Features and Linear Classification}

Social media platforms like Twitter are rich sources of emotional expression. In this homework, you'll build classifiers to detect emotions in tweets using linear classification with cross-entropy loss and softmax activation. Consider the following dataset of 6 tweets, each labeled with exactly one emotion: \textbf{joy} (J), \textbf{anger} (A), or \textbf{fear} (F):

\begin{table}[h]
\centering
\lstDeleteShortInline|
\begin{tabular}{|l|l|l|}
\hline
Tweet & Emotion & One-hot encoding \\
\hline
``amazing day" & Joy & [1, 0, 0] \\
``scared of spiders" & Fear & [0, 0, 1] \\
``love this" & Joy & [1, 0, 0] \\
``so angry" & Anger & [0, 1, 0] \\
``so so worried about tomorrow" & Fear & [0, 0, 1] \\
``hate waiting" & Anger & [0, 1, 0] \\
\hline
\end{tabular}
\lstMakeShortInline|
\end{table}

Each tweet $x$ is mapped to a feature vector $f(x)$ using \textbf{bag-of-words representation} (you can think of it as word counts). Machine learning models don't directly take a string as input; instead, as you have learned, they work with matrices (like NumPy arrays and tensors). One way to convert strings to tensors is “bag-of-words” features, which treat input texts as a "bag of words" and represents them using a vector with the same length as our vocabulary. The resulting representation is a vector, where each position corresponds to a word in the vocabulary, and each value represents how many times that word appears in the input text. If we have a very large vocabulary, we would end up with a very sparse vector with many 0's.\\

For this problem, let's assume our vocabulary consists of all unique words in the tweets above: \textbf{\{about, amazing, angry, day, hate, love, of, scared, so, spiders, this, tomorrow, waiting, worried\}}.\\

To build our multi-class classifier, we use a weight matrix $\mathbf{W} \in \mathbb{R}^{d \times 3}$ where $d$ is the feature dimension (vocabulary size)
and 3 is the number of classes. The model computes logits as $\mathbf{z} = \mathbf{W}^T f(x)$ and applies softmax to convert
logits into probability values that sum to 1. Let $\mathbf{p}$ be the 3-dimensional vector of output probabilities, and $p_k$ be
the probability that the $k$th element of $p$ has a true label of 1. The softmax probabilities are computed in the following way:
$$p_k = \frac{e^{z_k}}{\sum_{j=1}^{3} e^{z_j}}$$

Then, given the softmax probabilities, the classifier will use argmax to choose the prediction class with the highest $p_k$. \\

The cross-entropy loss measures the difference between a model's predicted probabilities and the true probability distribution of the data. For a single example (with three output classes), it can be computed as follows: $$\text{Loss}_{\text{CE}}(x, \mathbf{y}, \mathbf{W}) = -\sum_{k=1}^{3} y_k \log p_k$$
where $\mathbf{y}$ is the one-hot encoded true label vector.

\begin{enumerate}

  \input{01-bag-of-words/01-feature-vector}

  \input{01-bag-of-words/02-softmax}

  \input{01-bag-of-words/03-cross-entropy-loss}

  \input{01-bag-of-words/04-gradient-analysis}

\end{enumerate}
